\section*{第 \ref{ch:g11:the-eye} 章——眼睛及其光感受器}

\begin{solution}{exo:g11:the-eye:1}
角膜、瞳孔（经虹膜）、晶状体、玻璃体，然后是
视网膜的透明层（神经节\omterm{def:g10:cells-common-unit:cell}{细胞}和双极\omterm{def:g10:cells-common-unit:cell}{细胞}），再到
后部的\omterm{def:g11:the-eye:photoreceptors}{光感受器}。
\end{solution}

\begin{solution}{exo:g11:the-eye:2}
\omterm{def:g11:the-eye:photoreceptors}{视杆细胞}约 1.2 亿个，灵敏度极高，只有一种色素，夜间视觉呈
灰色。\omterm{def:g11:the-eye:photoreceptors}{视锥细胞}约 600 万个，需要强光，有三种色素，
负责白天视觉和色觉，在中央凹最为敏锐。
\end{solution}

\begin{solution}{exo:g11:the-eye:3}
中央凹是视网膜中央的凹陷，仅由排列紧密的
\omterm{def:g11:the-eye:photoreceptors}{视锥细胞}构成，视觉最为敏锐。盲点是
视神经穿出之处，没有\omterm{def:g11:the-eye:photoreceptors}{光感受器}。
\end{solution}

\begin{solution}{exo:g11:the-eye:4}
S \omterm{def:g11:the-eye:photoreceptors}{视锥细胞}约占 45\%，\omterm{def:g11:the-eye:photoreceptors}{视杆细胞}约占 60\%，M \omterm{def:g11:the-eye:photoreceptors}{视锥细胞}约占 20\%，L
\omterm{def:g11:the-eye:photoreceptors}{视锥细胞}约占 10\%。
\end{solution}

\begin{solution}{exo:g11:the-eye:5}
位于 X \omterm{def:g11:cell-cycle-mitosis:chromosome}{染色体}上；该病为 X 连锁\omterm{def:g11:genes-and-disease:genetic}{隐性}遗传，每个携带该\omterm{def:g10:universal-dna:gene}{等位基因}的
男性都会表现出来，女性只有两个拷贝都携带时才会表现，因此
在男性中远比女性常见。
\end{solution}

\begin{solution}{exo:g11:the-eye:6}
夜间只有\omterm{def:g11:the-eye:photoreceptors}{视杆细胞}起作用，且只含一种色素：它们
只能反映亮度而非波长，因此无色觉。中央凹只有
\omterm{def:g11:the-eye:photoreceptors}{视锥细胞}，灵敏度不足以感知这颗恒星；斜视时，恒星的
图像落在\omterm{def:g11:the-eye:photoreceptors}{视杆细胞}上，由它们探测到。
\end{solution}

\begin{solution}{exo:g11:the-eye:7}
颜色相同：两种情况下 S:M:L 之比均为 0:0.6:1。
亮度不同：光 2 使每个\omterm{def:g11:the-eye:photoreceptors}{视锥细胞}的兴奋减半。颜色由
比值决定，亮度由总量决定。
\end{solution}

\begin{solution}{exo:g11:the-eye:8}
儿子的 X \omterm{def:g11:cell-cycle-mitosis:chromosome}{染色体}来自母亲，全部色觉正常。女儿
从父亲获得 X \omterm{def:g11:cell-cycle-mitosis:chromosome}{染色体}，全部为携带者但色觉正常。每位携带者女儿的
儿子患色盲的概率为二分之一。
\end{solution}

\begin{solution}{exo:g11:the-eye:9}
男性只需一条\omterm{def:g11:mutations:mutation}{突变} X \omterm{def:g11:cell-cycle-mitosis:chromosome}{染色体}，女性需要两条。若该\omterm{def:g10:universal-dna:gene}{等位基因}的频率为
0.08，则女性携带两条的概率为 $0.08^2 = 0.0064$，约
0.6\%——比 8\% 少十六倍。
\end{solution}

\begin{solution}{exo:g11:the-eye:10}
差异随自复制事件发生以来的时间而积累：L 与 M 之间相差 4\%，
说明是近期的复制；二者与 S 之间相差 57\%，说明是
久远得多的复制。L/M 复制是最近发生的事件，S 的分化
则早得多。
\end{solution}

\begin{solution}{exo:g11:the-eye:11}
旧大陆灵长类的祖先复制了其单一的 X 连锁
\omterm{def:g11:the-eye:photoreceptors}{视蛋白}\omterm{def:g10:universal-dna:gene}{基因}；其中一份拷贝的吸收峰移向较长波长，产生了
第三种\omterm{def:g11:the-eye:photoreceptors}{视锥细胞}类型，并开辟了色觉的新维度。这一变化被保留下来，可能
因为它有助于在叶丛中辨认成熟的果实和嫩叶。
\end{solution}

\begin{solution}{exo:g11:the-eye:12}
蓝色与黄色（以及蓝绿色与绿色）。S \omterm{def:g10:universal-dna:gene}{基因}位于一条
普通\omterm{def:g11:cell-cycle-mitosis:chromosome}{染色体}上，因此无论男女，两个\omterm{def:g10:universal-dna:gene}{等位基因}都必须发生\omterm{def:g11:mutations:mutation}{突变}——
这种情况罕见，且与性别无关。
\end{solution}

\begin{solution}{exo:g11:the-eye:13}
$\num{2.5e-6}/\num{17e-3} \approx \num{1.5e-4}\,\unit{rad}$，约
半角分。在 \qty{10}{m} 处：$\num{1.5e-4} \times 10 =
\qty{1.5}{mm}$。
\end{solution}

\begin{solution}{exo:g11:the-eye:14}
男性只有一条 X \omterm{def:g11:cell-cycle-mitosis:chromosome}{染色体}，因而只有一个\omterm{def:g10:universal-dna:gene}{等位基因}，即只有一种
M/L 型\omterm{def:g11:the-eye:photoreceptors}{视锥细胞}，加上 S 型，共两种\omterm{def:g11:the-eye:photoreceptors}{视锥细胞}。携带两个不同\omterm{def:g10:universal-dna:gene}{等位基因}
的女性，部分\omterm{def:g11:the-eye:photoreceptors}{视锥细胞}表达其中一个，其余表达另一个，共三种，因此为
三色视觉；携带两个相同\omterm{def:g10:universal-dna:gene}{等位基因}的女性则与男性相同，只有两种。
\end{solution}

\begin{solution}{exo:g11:the-eye:15}
神经节\omterm{def:g10:cells-common-unit:cell}{细胞}和视神经完好，因此施加于它们的电脉冲
会作为信号传至视觉脑区。图像质量受限于电极数量
（仅数百个，而视神经纤维约一百万根、\omterm{def:g11:the-eye:photoreceptors}{光感受器}
数百万个），以及视网膜自身对图像的分类处理
功能的丧失。
\end{solution}

\begin{solution}{pb:g11:the-eye:1}
\textbf{1.} \qty{470}{nm}：S 40，M 35，L 18。\qty{530}{nm}：S 0，M 100，
L 80。\qty{590}{nm}：S 0，M 40，L 90。\qty{650}{nm}：S 0，M 2，L 12。

\textbf{2.} \qty{470}{nm} 同时激发三种\omterm{def:g11:the-eye:photoreceptors}{视锥细胞}；\qty{650}{nm} 实际上
只激发 L 型\omterm{def:g11:the-eye:photoreceptors}{视锥细胞}。

\textbf{3.} 仅考虑 S 与 M 时：\qty{590}{nm} 给出 (0, 40)，而
\qty{650}{nm} 给出 (0, 2)；\qty{530}{nm} 给出 (0, 100)。
\qty{590}{nm}（橙色）与 \qty{530}{nm}（绿色）两种光仅在
M 的兴奋强度上不同，类似一暗一亮的两种绿色，是难以区分的一对；
\qty{650}{nm} 看起来像非常暗的光。

\textbf{4.} 仅考虑 S 与 L 时：\qty{530}{nm} 给出 (0, 80)，
\qty{590}{nm} 给出 (0, 90)——二者几乎相同：绿色与橙色
易被混淆。

\textbf{5.} 二者仍能探测红光，依靠其保留的\omterm{def:g11:the-eye:photoreceptors}{视锥细胞}类型；
失去的是 M 与 L 响应之间的比较，
而这一比较正是区分红色与绿色的关键。

\textbf{6.} $X^NX^c$：她的父亲把唯一的 X 传给了她，这条 X 携带\omterm{def:g11:mutations:mutation}{突变}\omterm{def:g10:universal-dna:gene}{等位基因}；
她色觉正常，说明另一条 X 是正常的。

\textbf{7.} 儿子：$X^NY$ 正常或 $X^cY$ 色盲，各占一半。
女儿：$X^NX^N$ 或 $X^NX^c$，色觉均正常，各占一半
（后者为携带者）。

\textbf{8.} 伊内斯为携带者的概率为 $1/2$。女儿若从双亲各获得一条\omterm{def:g11:mutations:mutation}{突变}的 X，
则为色盲：从父亲处必然获得 $X^c$，
从伊内斯处获得的概率为 $1/2 \times 1/2$：$1/4$。

\textbf{9.} 色盲女儿有两条\omterm{def:g11:mutations:mutation}{突变}的 X \omterm{def:g11:cell-cycle-mitosis:chromosome}{染色体}，分别来自父母各一方；
父亲把唯一的 X 传给每个女儿，
因此他的 X 是\omterm{def:g11:mutations:mutation}{突变}的，他本人为色盲。

\textbf{10.} 是的：她的儿子从她那里获得 X，这条 X 携带
\omterm{def:g11:mutations:mutation}{突变}\omterm{def:g10:universal-dna:gene}{等位基因}；她色觉正常，因此她的\omterm{def:g11:enzymes-and-phenotype:phenotype}{基因型}为 $X^NX^c$。

\textbf{11.} 视紫红质最先分出；随后 S 从 L/M 线中分出；
最后 L 与 M 分开：（视紫红质，（S，（L，M）））。

\textbf{12.} 差异 4\% 对应 3500 万年；差异 57\% 则需
约 $35 \times 57/4 \approx 500$ 百万年——与
\omterm{def:g10:common-ancestry:bodyplan}{脊椎动物}的早期起源相符（二者关系在差异较大时并非严格成正比，
因此这里只是数量级估计）。

\textbf{13.} 两个几乎相同的相邻\omterm{def:g10:universal-dna:gene}{基因}可能在两条 X \omterm{def:g11:cell-cycle-mitosis:chromosome}{染色体}配对时错位；
随后交换产生
一条含三个拷贝的 X 和一条只含单个拷贝的 X。获得单拷贝 X 的男性
只有一种 M/L 型\omterm{def:g11:the-eye:photoreceptors}{视锥细胞}，
即红绿色盲。

\textbf{14.} 配对依赖序列相似性：两个相邻的、序列相同率为 96\%
的\omterm{def:g10:universal-dna:gene}{基因}可能错位配对，即\omterm{def:g10:universal-dna:gene}{基因} 1 与\omterm{def:g10:universal-dna:gene}{基因} 2 配对，
并发生不等交换。

\textbf{15.} 该变化只发生过一次，发生于旧大陆猴类的共同祖先中，
在其与其他哺乳动物及新大陆猴分开之后，
约 3500 万年前，
并被其所有后代遗传。

\textbf{16.} $\num{2.5e-6}/\num{17e-3} \approx \num{1.5e-4}$ 弧度，
约为 $0.5'$。

\textbf{17.} $\num{1.5e-4} \times 6 \approx \qty{0.9}{mm}$：比
\qty{1.7}{mm} 的笔画更细，因此正常视力能轻松辨认该行
（实际上光学模糊使极限接近
$1'$）。

\textbf{18.} 至少粗十倍（100 个\omterm{def:g11:the-eye:photoreceptors}{视杆细胞}共享一个信号，
约 10 乘 10）：分辨角为数角分；作为代价，一百个
感受器汇集其光能，从而能探测暗得多的光源。

\textbf{19.} 密度加倍，间距为原来的 $\sqrt{2}$ 分之一，并
眼球增大则焦距变长：分辨率约提高为两倍。

\textbf{20.} S、M 和 L 三种\omterm{def:g11:the-eye:photoreceptors}{视锥细胞}，各自对光谱中远离其峰值的另一端几乎不敏感；
X 连锁\omterm{def:g11:the-eye:photoreceptors}{视蛋白}\omterm{def:g10:universal-dna:gene}{基因}的一次串联重复，
发生于 3500 万年前，使灵长类获得了第三种\omterm{def:g11:the-eye:photoreceptors}{视锥细胞}；中央凹的分辨
率约为半角分。
\end{solution}
